\chapter{KIỂM THỬ THỰC NGHIỆM VÀ ĐỐI KHÁNG}

\section{Chiến lược kiểm thử đa tầng (Testing Pyramid)}

Nhằm bảo đảm hệ thống MarketFlow AI vận hành an toàn và ổn định trong môi trường sản xuất thực tế, dự án thiết lập một chiến lược kiểm thử tự động đa tầng theo mô hình Kim tự tháp kiểm thử (Testing Pyramid) với tổng cộng \textbf{151 ca kiểm thử tự động} và \textbf{8 kịch bản đánh giá định lượng Prompt}:
\begin{enumerate}[noitemsep,topsep=2pt]
 \item \textbf{Tầng Kiểm thử Backend API \& Máy chủ FastMCP (\texttt{backend/tests/})}: 104 ca kiểm thử tự động trên 16 tệp kiểm thử bằng \texttt{pytest} và \texttt{fastapi.testclient}, chạy trên cơ sở dữ liệu cô lập tự động làm sạch sau mỗi ca test.
 \item \textbf{Tầng Kiểm thử Dịch vụ Điều phối AI (\texttt{ai-runner/tests/})}: 18 ca kiểm thử đơn vị và bất đồng bộ (\texttt{pytest-asyncio}) kiểm định hàng đợi FIFO \texttt{RUNNER\_CONCURRENCY=1}, chuyển đổi luồng NDJSON sang SSE, xử lý timeout và xoay vòng đa mô hình dự phòng.
 \item \textbf{Tầng Kiểm thử Đầu-Cuối trên Trình duyệt thực (\texttt{frontend/tests/e2e/})}: 29 ca kiểm thử tự động bằng \textbf{Playwright} trên trình duyệt Chromium thực tế, kiểm tra trọn vẹn luồng giao diện từ đăng nhập, dashboard, chiến dịch, duyệt bài, lịch đăng, nhập CSV đến khung chat AI.
\end{enumerate}

\section{Danh mục các bộ kiểm thử tự động (151 ca kiểm thử)}

Bảng~\ref{tab:test-suites} trình bày chi tiết từng tệp kiểm thử tự động trong kho mã nguồn thực tế của dự án cùng số lượng ca kiểm thử và kết quả chạy thực nghiệm.

\begin{small}
\begin{longtable}{|L{4.3cm}|C{1.4cm}|L{6.1cm}|C{1.9cm}|}
\caption{Danh mục 27 tệp kiểm thử tự động trong MarketFlow AI (151 ca kiểm thử)}\label{tab:test-suites}\\
\hline
\thead{Tệp kiểm thử} & \thead{Số ca} & \thead{Trọng tâm kiểm định kỹ thuật} & \thead{Tỷ lệ Pass}\\
\hline
\endfirsthead
\hline
\thead{Tệp kiểm thử} & \thead{Số ca} & \thead{Trọng tâm kiểm định kỹ thuật} & \thead{Tỷ lệ Pass}\\
\hline
\endhead
\hline
\endfoot
\hline
\endlastfoot
\textbf{Nhóm A (Backend)} & \textbf{104} & \textbf{Bộ kiểm thử Backend API \& FastMCP (\texttt{backend/tests/} -- 16 tệp)} & \textbf{100\%}\\
\hline
\texttt{test\_auth.py} & 12 & Đăng nhập hợp lệ/sai mật khẩu (401), giới hạn tần suất (429), tài khoản bị khóa, tự đăng ký, thu hồi token khi đăng xuất (\texttt{token\_version}), đổi mật khẩu, lấy thông tin \texttt{/me}. & 100\% (12/12)\\
\hline
\texttt{test\_users.py} & 6 & Phân quyền \texttt{MANAGER} tạo/sửa/khóa tài khoản; chặn \texttt{MARKETER} truy cập \texttt{/users} (403); chống tự khóa chính mình (\texttt{CANNOT\_MODIFY\_SELF}). & 100\% (6/6)\\
\hline
\texttt{test\_channels.py} & 5 & Danh sách 7 kênh, tạo kênh mới bởi \texttt{MANAGER}, chặn xóa kênh đang dùng (\texttt{CHANNEL\_IN\_USE} 409), chặn \texttt{MARKETER} sửa đổi kênh (403). & 100\% (5/5)\\
\hline
\texttt{test\_campaigns.py} & 7 & Tạo chiến dịch, kiểm tra \texttt{budget >= 0}, \texttt{start\_date <= end\_date}, phân bổ kênh, kiểm soát quyền sở hữu \texttt{owner\_id} (403), chặn xóa khi đã có dữ liệu (409). & 100\% (7/7)\\
\hline
\texttt{test\_contents.py} & 9 & Chu trình HITL: tạo \texttt{DRAFT}, quét từ khóa cấm khi \texttt{/submit}, chặn tự duyệt \texttt{CANNOT\_REVIEW\_OWN}, duyệt/từ chối, đánh dấu xuất bản, khóa xóa \texttt{CONTENT\_LOCKED} (409). & 100\% (9/9)\\
\hline
\texttt{test\_schedules.py} & 7 & Lập lịch cho bài viết hợp lệ, kiểm tra khung ngày chiến dịch, kiểm soát quyền sở hữu, hoàn tất (\texttt{/done}) và hủy lịch đăng. & 100\% (7/7)\\
\hline
\texttt{test\_metrics.py} & 9 & Kiểm tra bất biến \texttt{views >= clicks >= conversions}, gộp theo ngày, kiểm soát quyền sở hữu, nhập CSV 2 bước (\texttt{import/preview} \& \texttt{import/commit}). & 100\% (9/9)\\
\hline
\texttt{test\_stats.py} & 6 & Kiểm định công thức CTR, CVR, CPC, CPA, CPM, chống chia cho 0, thống kê theo chuỗi thời gian/kênh/chiến dịch và tỷ lệ giải ngân (\texttt{used\_pct}). & 100\% (6/6)\\
\hline
\texttt{test\_brand.py} & 1 & Đọc và cập nhật trọn vẹn hồ sơ thương hiệu (\texttt{brand\_profile}) cùng danh sách từ khóa cấm (\texttt{banned\_keywords}). & 100\% (1/1)\\
\hline
\texttt{test\_notifications.py} & 1 & Kiểm thử chu trình thông báo đầu-cuối: danh sách thông báo, đếm chưa đọc, đánh dấu đã đọc từng mục và toàn bộ. & 100\% (1/1)\\
\hline
\texttt{test\_health.py} & 3 & Kiểm tra điểm cuối \texttt{/api/v1/health} khi cục bộ, khi \texttt{ai-runner} sẵn sàng và khi mất kết nối (\texttt{degraded}). & 100\% (3/3)\\
\hline
\texttt{test\_assistant\_mcp.py} & 6 & Xác thực \texttt{Run-Token} JWT tại \texttt{/mcp}, từ chối token hết hạn/sai kiểu \texttt{typ}/đã kết thúc lượt, thực thi công cụ qua giao thức JSON-RPC MCP. & 100\% (6/6)\\
\hline
\texttt{test\_assistant\_tools.py} & 4 & Kiểm thử trực tiếp 9 công cụ đọc và 5 công cụ đề xuất \texttt{propose\_*} tạo bản ghi \texttt{ai\_actions} ở trạng thái \texttt{PROPOSED} kèm chống trùng lặp. & 100\% (4/4)\\
\hline
\texttt{test\_assistant\_}\allowbreak\texttt{actions.py} & 9 & Xác nhận (\texttt{/confirm}) và từ chối (\texttt{/reject}) cho các loại \texttt{ai\_action}; ghi đè tham số, chặn xác nhận lặp lại (409) và kiểm tra quyền sở hữu. & 100\% (9/9)\\
\hline
\texttt{test\_assistant\_}\allowbreak\texttt{stream.py} & 14 & Luồng SSE hội thoại AI, quản lý phiên, dừng luồng (\texttt{/stop}), tạo lại (\texttt{/regenerate}), giữ khóa phiên và báo lỗi trung thực khi hết hạn mức. & 100\% (14/14)\\
\hline
\texttt{test\_assistant\_}\allowbreak\texttt{context.py} & 5 & Tạo snapshot ngữ cảnh ở lượt đầu không bịa số liệu, kèm chỉ số thực theo kênh, thu gọn ở các lượt sau, tách phiên mới khi đổi chiến dịch. & 100\% (5/5)\\
\hline
\textbf{Nhóm B (AI-Runner)} & \textbf{18} & \textbf{Bộ kiểm thử Dịch vụ AI Runner (\texttt{ai-runner/tests/} -- 2 tệp)} & \textbf{100\%}\\
\hline
\texttt{test\_api.py} \& \texttt{test\_runner.py} & 18 & Xác thực khóa nội bộ \texttt{X-Runner-Token}, hàng đợi tuần tự \texttt{Semaphore(1)}, biên dịch sự kiện NDJSON sang SSE, tự động chuyển đổi mô hình dự phòng khi lỗi/timeout. & 100\% (18/18)\\
\hline
\textbf{Nhóm C (E2E UI)} & \textbf{29} & \textbf{Bộ kiểm thử Đầu-Cuối Playwright E2E (\texttt{frontend/tests/e2e/} -- 9 tệp)} & \textbf{100\%}\\
\hline
9 tệp \texttt{*.spec.ts} (\texttt{auth}, \texttt{dashboard}, \texttt{campaigns}, \texttt{contents\_reviews}, \texttt{calendar}, \texttt{metrics}, \texttt{assistant}, \texttt{settings}, \texttt{visual\_audit}) & 29 & Kiểm thử tự động trên trình duyệt Chromium: đăng nhập/đăng ký/đăng xuất, xem biểu đồ KPI, tạo chiến dịch, gửi duyệt \& phê duyệt bài viết, lập lịch, nhập CSV, mở ngăn kéo AI và kiểm tra giao diện. & 100\% (29/29)\\
\hline
\textbf{Tổng cộng toàn dự án} & \textbf{151} & \textbf{Bao phủ toàn diện Backend, FastMCP, AI-Runner và Giao diện E2E} & \textbf{100\% Pass}\\
\hline
\end{longtable}
\end{small}

\section{Kiểm thử đối kháng bảo mật và quy trình nghiệp vụ (Adversarial Testing)}

Nhằm bảo đảm độ tin cậy khách quan của kết quả nghiệm thu, nhóm đã thiết kế các ca kiểm thử đối kháng trực tiếp trong bộ test để giả lập các hành vi vượt quyền và gian lận quy trình nghiệp vụ (Bảng~\ref{tab:adversarial-matrix}).

\begin{table}[!htb]
\centering
\scriptsize
\renewcommand{\arraystretch}{0.94}
\caption{Ma trận kiểm thử đối kháng và cơ chế phòng thủ của hệ thống}\label{tab:adversarial-matrix}
\begin{tabularx}{\textwidth}{|L{3.1cm}|Y|L{3.2cm}|C{1.9cm}|}
\hline
\thead{Hình thức tấn công} & \thead{Kịch bản giả lập đối kháng} & \thead{Cơ chế phòng thủ \& Mã phản hồi} & \thead{Kết quả}\\
\hline
Tự phê duyệt bài viết của chính mình (Self-Approval) & Người dùng có vai trò \texttt{MANAGER} tự tạo bài viết rồi gọi \texttt{POST /contents/\{id\}/approve} trên bài viết đó. & Hàm \texttt{approve\_content} kiểm tra \texttt{content.created\_by == current\_user.id} và chặn với mã \texttt{CANNOT\_REVIEW\_OWN}. & Kháng cự 100\%\\
\hline
Leo quyền ngang giữa các Marketer (IDOR) & \texttt{marketer2} gửi yêu cầu sửa chiến dịch, thêm nội dung hoặc xác nhận thẻ \texttt{ai\_action} của \texttt{marketer1}. & Kiểm tra \texttt{campaign.owner\_id != current\_user.id}, từ chối ngay với mã \texttt{PERMISSION\_DENIED} (HTTP 403). & Kháng cự 100\%\\
\hline
Xóa bài viết đã duyệt để xóa dấu vết & Tác giả hoặc quản lý gọi \texttt{DELETE /contents/\{id\}} đối với bài viết đang ở trạng thái \texttt{APPROVED} hoặc \texttt{PUBLISHED}. & Chặn xóa với mã lỗi \texttt{CONTENT\_LOCKED} (HTTP 409 Conflict) nhằm bảo vệ lịch sử xuất bản. & Kháng cự 100\%\\
\hline
Gửi duyệt nội dung chứa từ khóa cấm & Tác giả đưa cụm từ thuộc danh mục từ cấm (ví dụ: ``cam kết 100\%'') vào bài viết và gọi \texttt{/submit}. & Hàm \texttt{scan\_banned\_keywords} phát hiện từ khóa cấm và đính kèm danh sách \texttt{banned\_matches} cảnh báo cho người duyệt. & Kháng cự 100\%\\
\hline
Giả mạo Token người dùng gọi vào \texttt{/mcp} & Sử dụng Access Token đăng nhập thông thường hoặc token hết hạn để gọi trực tiếp máy chủ công cụ FastMCP \texttt{/mcp}. & Kiểm tra trường \texttt{typ == "assistant\_run"}; từ chối mọi token không phải \texttt{Run-Token} hợp lệ với HTTP 401. & Kháng cự 100\%\\
\hline
Nhập số liệu phễu tiếp thị phi lý & Gửi dữ liệu chỉ số có \texttt{clicks > views} hoặc \texttt{conversions > clicks} qua API hoặc tệp CSV. & Bộ kiểm định Pydantic v2 và trình phân tích CSV báo lỗi vi phạm bất biến số học (HTTP 422). & Kháng cự 100\%\\
\hline
\end{tabularx}
\end{table}

\section{Kết quả thực thi kiểm thử tổng hợp}

Bảng~\ref{tab:test-breakdown} tổng hợp kết quả thực thi thực tế của toàn bộ các bộ kiểm thử tự động trong dự án.

\begin{table}[!htb]
\centering
\footnotesize
\renewcommand{\arraystretch}{0.98}
\caption{Tổng hợp kết quả thực thi 151 ca kiểm thử tự động theo phân hệ}
\label{tab:test-breakdown}
\begin{tabularx}{\textwidth}{|L{4.8cm}|C{1.3cm}|C{1.3cm}|C{1.3cm}|Y|}
\hline
\thead{Phân hệ kiểm định} & \thead{Số ca} & \thead{Đạt} & \thead{Lỗi} & \thead{Thời gian \& Đánh giá kỹ thuật}\\
\hline
Backend Core CRUD, RBAC \& HITL (11 tệp test) & 66 & 66 & 0 & Xác thực JWT, \texttt{owner\_id}, \texttt{CANNOT\_REVIEW\_OWN}, CSV 2 bước, tính toán KPI.\\
\hline
Backend FastMCP \& Trợ lý AI (5 tệp test) & 38 & 38 & 0 & Xác thực \texttt{Run-Token}, 14 công cụ MCP, luồng SSE, xác nhận thẻ \texttt{ai\_action}.\\
\hline
Dịch vụ AI Runner cô lập (2 tệp test) & 18 & 18 & 0 & Thực thi hoàn tất trong \texttt{3,95s}; kiểm chứng hàng đợi tuần tự và xoay vòng đa mô hình.\\
\hline
Giao diện Đầu-Cuối Playwright E2E (9 tệp spec) & 29 & 29 & 0 & Kiểm chứng trực tiếp trên trình duyệt Chromium; kết hợp Ruff/TSC/ESLint 0 lỗi.\\
\hline
\textbf{Tổng cộng toàn hệ thống} & \textbf{151} & \textbf{151} & \textbf{0} & \textbf{Đạt 100,0\% (0 Failures, 0 Warnings)}\\
\hline
\end{tabularx}
\end{table}

\section{Kết luận chương}

Chương 6 đã cung cấp bằng chứng thực nghiệm đầy đủ về chất lượng mã nguồn của MarketFlow AI thông qua 151 ca kiểm thử tự động (104 Backend Pytest + 18 AI-Runner Pytest + 29 Playwright E2E) đạt tỷ lệ vượt qua 100\%. Chương tiếp theo sẽ trình bày kết quả đo đạc P0 Spike và đánh giá định lượng 4 phiên bản System Prompt của Trợ lý AI.
